\section{噪声感知人群的聚类分析}

探索性因子分析将居民的噪声影响感知压缩为四个维度，但不同个体在这四个维度上的组合方式存在差异。为识别具有差异化噪声感知特征的居民群体，本节在因子得分基础上进行聚类分析。

\subsection{聚类方法与最优类别数确定}

本研究以四个因子得分 $F_1\sim F_4$ 为输入特征，采用 KMeans++ 算法对467个样本进行聚类。KMeans++ 通过优化初始聚类中心的选择，降低传统 KMeans 对初始中心点的敏感性。为确定最优类别数，本文在 $K=2,3,\dots,8$ 的范围内分别计算轮廓系数与簇内误差平方和，并以轮廓系数最大作为选择准则。

\begin{table}[H]
\centering
\caption{不同类别数下的聚类评价指标}
\rowcolors{2}{tablight}{white}
\begin{tabular}{M{4cm}M{5cm}M{4cm}}
\bluetoprule
\rowcolor{tabhead}
\thead{类别数 $K$} & \thead{簇内误差平方和（SSE）} & \thead{轮廓系数} \\
\hline
2 & 1535.03 & 0.2046 \\
3 & 1262.56 & 0.2377 \\
4 & 1042.30 & 0.2549 \\
5 & 874.99 & 0.2847 \\
6 & 769.56 & 0.2623 \\
7 & 680.61 & 0.2773 \\
8 & 620.58 & 0.2836 \\
\bluebottomrule
\end{tabular}
\end{table}

轮廓系数衡量样本与同簇内其他样本的相似程度相对其与相邻簇样本相似程度的优势，取值介于 $-1$ 与 $1$ 之间，越接近 $1$ 表示聚类越紧凑、簇间分离度越高。经比较，$K=5$ 时轮廓系数最大，故将居民划分为五类。

\subsection{聚类结果与人群画像}

聚类结果通过降维散点图与因子均值雷达图共同呈现，前者刻画样本在低维空间中的分离程度，后者刻画各类人群在四个因子上的得分结构，如图\ref{fig:cluster_profile}所示。

\begin{figure}[H]
\centering
\begin{subfigure}[b]{0.48\textwidth}
  \centering
  \includegraphics[width=\textwidth]{cluster_3d.png}
  \caption{聚类结果的主成分三维散点图}
\end{subfigure}
\hfill
\begin{subfigure}[b]{0.48\textwidth}
  \centering
  \includegraphics[width=\textwidth]{cluster_radar.png}
  \caption{各人群类型的因子均值雷达图}
\end{subfigure}
\caption{噪声感知人群聚类结果与人群画像}
\label{fig:cluster_profile}
\end{figure}

图\ref{fig:cluster_profile}(a)将四个因子得分经主成分分析降至三维，并按所属类别着色绘制散点图，用于直观呈现聚类结果的样本分布。图中五类样本在低维空间中形成相对分离的团簇，簇间边界较为清晰、簇内分布相对紧凑，说明以 $K=5$ 划分具有合理性。从各类规模看，样本量差异明显：其中一类包含211个样本，约占有效样本的45.2\%；另有三个规模相近的类别分别包含73、69和68个样本；最小的一类为46个样本。

图\ref{fig:cluster_profile}(b)进一步给出上述五类人群在四个噪声影响因子上的均值轮廓。五类之间的差异\textbf{并不体现为困扰程度的高低，而在于困扰来源的结构}：\textbf{一般型}（规模最大）四个维度均贴近总体均值，尚未形成突出的噪声困扰；\textbf{低敏感低困扰型}（规模最小）四个因子得分同时为负，困扰程度最低。其余三类各有侧重：\textbf{身心情绪高敏感型}在\textbf{身心健康与认知干扰}与\textbf{情绪与社会关系影响}两个维度上同时外凸，其困扰主要来自自身对噪声的生理与情绪反应；\textbf{典型环境噪声暴露型}以\textbf{典型环境噪声源干扰}为唯一突出维度，困扰集中于交通、施工、公共活动等公共性噪声源；\textbf{生活场景噪声困扰型}以\textbf{生活场景噪声源干扰}最为突出，困扰更多来自居住环境中贴近个体的人为噪声。

为刻画各类人群在原始维度上的分离程度，本文对四个因子分别进行单因素方差分析，结果表明四个因子在五类之间的组间差异均达到极显著水平（$P<0.001$）。需要说明的是，这四个因子同时是聚类的输入特征，该结果反映的是聚类对原始因子信息的有效划分，而非聚类结构的外部效度证据。根据各簇在四个因子上的得分特征，将五类人群命名如下：

\begin{table}[H]
\centering
\caption{五类噪声感知人群的因子得分特征}
\rowcolors{2}{tablight}{white}
\begin{tabular}{M{3.5cm}M{8cm}}
\bluetoprule
\rowcolor{tabhead}
\thead{人群类型} & \thead{因子得分特征} \\
\hline
低敏感低困扰型 & $F_1\sim F_4$ 均低于平均水平 \\
一般型 & 各因子接近总体均值，为最大群体 \\
生活场景噪声困扰型 & $F_4$ 生活场景噪声源干扰突出 \\
典型环境噪声暴露型 & $F_3$ 典型环境噪声源干扰突出 \\
身心情绪高敏感型 & $F_1$ 与 $F_2$ 同时偏高 \\
\bluebottomrule
\end{tabular}
\end{table}

聚类结果说明，居民的噪声影响感知并非均匀分布，而是呈现出若干可区分的类型。其中身心情绪高敏感型人群在身心健康与情绪社会关系两个维度上均处于高位，是噪声治理中需要重点关注的对象。

